\documentclass[conference]{IEEEtran}
\IEEEoverridecommandlockouts

\IfFileExists{ptmr7t.tfm}{}{%
  \typeout{** Times (ptmr7t.tfm) not found: loading Latin Modern. Install texlive-fonts-recommended for IEEE Times.}%
  \usepackage{lmodern}%
}

\usepackage[T1]{fontenc}
\usepackage{amsmath,amssymb,amsfonts}
\usepackage{graphicx}
\usepackage{textcomp}
\usepackage{url}

\newcommand{\PaperTitle}{Nueva Perspectiva del Nuevo Juego de Desarrollo de Nuevos Productos}

\begin{document}
\renewcommand{\abstractname}{Resumen}
\renewcommand{\IEEEkeywordsname}{Términos clave}
\renewcommand{\refname}{Referencias}

\title{\PaperTitle}

\author{\IEEEauthorblockN{Alejandro Castro López}
\IEEEauthorblockA{\textit{MNA} \\
\textit{ITESM}\\
Matrícula A00806170\\
a00806170@tec.mx}
}

\maketitle

\begin{abstract}
Reseña en español del artículo clásico de Takeuchi y Nonaka (1986) sobre desarrollo holístico de nuevos productos, publicado en \emph{Harvard Business Review}. Se sintetizan el contraste entre el modelo secuencial tipo relevo y el enfoque tipo rugby, las seis características observadas en empresas japonesas y su relevancia para velocidad y flexibilidad en mercados competitivos.
\end{abstract}

\begin{IEEEkeywords}
desarrollo de nuevos productos, gestión de proyectos, equipos multifuncionales, innovación, Takeuchi, Nonaka.
\end{IEEEkeywords}

\section{Presentación}

Hirotaka Takeuchi es profesor en la unidad de estrategia de Harvard Business School; Ikujiro Nonaka fue profesor emérito en la Universidad Hitotsubashi (Tokio) y es una figura central en la teoría de la creación de conocimiento organizacional \cite{takeuchiNonaka1986nndp}.

En enero de 1986 publicaron en \emph{Harvard Business Review} el artículo \emph{The New New Product Development Game} (en español, \emph{El nuevo juego de desarrollo de nuevos productos}) \cite{takeuchiNonaka1986nndp}. El texto analiza cómo varias empresas japonesas (entre ellas Fuji-Xerox, Canon, Honda y NEC) lograron acortar plazos y mejorar resultados en el desarrollo de nuevos productos durante la década de 1970 y principios de la de 1980, en contraste con el modelo occidental predominante de fases secuenciales y especialización estrecha.

\section{Resumen}

Takeuchi y Nonaka sostienen que las reglas del juego en el desarrollo de nuevos productos están cambiando: ya no basta con calidad, bajo costo y diferenciación; hace falta \textbf{velocidad} y \textbf{flexibilidad} para competir \cite{takeuchiNonaka1986nndp}. Proponen sustituir la metáfora del \textbf{relevo}---en la que cada función recibe el trabajo, lo termina y lo pasa al siguiente eslabón---por la del \textbf{rugby}, donde el equipo avanza junto hacia la meta y el ``balón'' circula dentro del grupo con movimiento continuo y coordinado \cite{takeuchiNonaka1986nndp}.

El núcleo del argumento es una \textbf{estrategia holística}: un equipo de desarrollo actúa como unidad con un objetivo común, en lugar de optimizar por departamentos de forma aislada. La alta dirección suele fijar una dirección estratégica amplia y metas exigentes, pero no entrega en general un concepto de producto cerrado ni un plan de trabajo detallado; el equipo dispone de margen para definir el concepto y la ejecución cotidiana---con un rol de la dirección más cercano al de un capitalista de riesgo que al de un controlador de tareas \cite{takeuchiNonaka1986nndp}.

A partir del estudio de casos japoneses, los autores identifican \textbf{seis características} que, en conjunto, configuran el nuevo juego:

\begin{enumerate}
    \item \textbf{Inestabilidad incorporada (\emph{built-in instability}):} proyectos estratégicamente importantes con objetivos ambiciosos y cierta tensión deliberada, que empujan a la creatividad sin microgestionar cada paso.
    \item \textbf{Equipos de proyecto autoorganizados:} grupos multifuncionales, motivados y con autonomía operativa; la autoorganización implica autonomía, superación de metas propias (\emph{self-transcendence}) y fertilización cruzada entre especialidades y estilos de pensamiento, no ausencia de dirección desde arriba.
    \item \textbf{Solapamiento de fases de desarrollo:} las etapas (por ejemplo, planificación, diseño, ingeniería, manufactura) se superponen en el tiempo en lugar de ejecutarse de forma estrictamente secuencial.
    \item \textbf{Multiaprendizaje (\emph{multilearning}):} aprendizaje simultáneo a nivel individual, de equipo, de la empresa y entre funciones, con experimentación iterativa y aceptación de decisiones revisadas cuando llega información nueva.
    \item \textbf{Control sutil:} diseño deliberado de incentivos, selección de integrantes, entorno de trabajo, contacto con clientes, evaluación por equipo y mecanismos que orientan el comportamiento deseado sin burocracia rígida.
    \item \textbf{Transferencia organizacional del aprendizaje:} disposición explícita a difundir lo aprendido en un proyecto al resto de la organización (por ejemplo, rotación, círculos de calidad, tiempo para innovación lateral).
\end{enumerate}

Los autores contrastan tres modos de coordinación temporal: el \textbf{sistema tipo A} (fases secuenciales, división del trabajo clara y evaluación individual), el \textbf{tipo B} (solapamiento entre fases adyacentes) y el \textbf{tipo C} (solapamiento amplio entre casi todas las fases). En B y C se espera trabajo con información parcial, retrabajo controlado y una \textbf{división del trabajo compartida}: los miembros pueden contribuir más allá de su especialidad formal---lo que exige perfiles en ``T'' y colaboración intensa \cite{takeuchiNonaka1986nndp}.

El artículo no es un manual de metodología ágil, pero sentó bases conceptuales ampliamente citadas en prácticas posteriores (incluido Scrum): equipos pequeños multifuncionales, iteración, entrega incremental de valor y aprendizaje continuo frente a un plan monolítico definido al inicio \cite{takeuchiNonaka1986nndp}. Su valor permanente está en explicar \emph{por qué} la velocidad competitiva exige repensar la arquitectura organizacional del desarrollo de productos, no solo optimizar cada fase por separado.

\section{Comentario crítico}

% Sección reservada para el autor.

\section{Conclusiones}

% Sección reservada para el autor.

\begin{thebibliography}{00}

\bibitem{takeuchiNonaka1986nndp}
H. Takeuchi and I. Nonaka, ``The New New Product Development Game,'' \emph{Harvard Bus. Rev.}, vol.~64, no.~1, pp.~137--146, Jan.--Feb. 1986. [En línea]. Disponible: \url{https://hbr.org/1986/01/the-new-new-product-development-game}

\end{thebibliography}

\end{document}
